\documentclass[12pt, a4paper]{article}

\usepackage[T1]{fontenc}
\usepackage[utf8]{inputenc}
\usepackage[british]{babel}
\usepackage{mathtools}
\usepackage{amsthm}
\usepackage{libertine}
\usepackage[libertine]{newtxmath}
\usepackage[mathscr]{euscript}
\usepackage{listings}
\usepackage{enumitem}
\usepackage{tikz-cd}
\usetikzlibrary{decorations.markings}
\usepackage{float}
\usepackage{geometry}
 \geometry{
 a4paper,
 total={170mm,257mm},
 left=20mm,
 top=30mm,
 }
\usepackage{fancyhdr}
\usepackage{hyperref}
\urlstyle{same}
\usepackage[noabbrev]{cleveref}
\usepackage{xcolor}
\usepackage{tcolorbox}
\tcbuselibrary{listings}

% Colors to match VS Code Lean 4 Light Theme
\definecolor{keywordcolor}{rgb}{0.0, 0.0, 1.0}   % blue
\definecolor{tacticcolor}{rgb}{0.0, 0.0, 1.0}    % blue
\definecolor{commentcolor}{rgb}{0.0, 0.5, 0.0}   % green
\definecolor{symbolcolor}{rgb}{0.0, 0.0, 0.0}    % black
\definecolor{sortcolor}{rgb}{0.0, 0.0, 1.0}      % blue
\definecolor{attributecolor}{rgb}{0.0, 0.0, 0.0} % black
% Light gray background color
\definecolor{codebg}{rgb}{0.97, 0.97, 0.97}

\def\lstlanguagefiles{lstlean.tex}

\newtcblisting{leancode}{
  listing only,
  listing options={
    language=lean,
    emph={sorry},
    emphstyle=\color{red},
    basicstyle=\ttfamily\footnotesize
  },
  colback=codebg,
  colframe=codebg,
  boxrule=0pt,
  arc=3mm,
  auto outer arc,
  top=10pt,
  bottom=10pt,
  left=5pt,
  right=5pt,
  width=\textwidth,
}

\theoremstyle{definition}
\newtheorem{innercustomthm}{Exercise}
\newenvironment{xca}[1] {\renewcommand\theinnercustomthm{#1}\innercustomthm}
  {\endinnercustomthm}

\pagestyle{fancy}
\fancyhf{}
\rhead{Lean Algebra Exercises{\renewcommand{\thefootnote}{\fnsymbol{footnote}}\footnotemark[2]}}
\chead{Nilradical}
\lhead{Pedro N\'{u}\~{n}ez}
% \rfoot{Page \thepage}

\title{Lean Algebra Exercise 6.4}
\author{Pedro N\'{u}\~{n}ez}

\setcounter{tocdepth}{1}
\sloppy
\makeatletter
\hypersetup{
  pdfauthor={\@author},
  pdftitle={\@title},
  colorlinks,
  linkcolor=[rgb]{0.2,0.2,0.6},
  citecolor=[rgb]{0.2,0.2,0.6},
  urlcolor=[rgb]{0.2,0.2,0.6}}
\makeatother

\begin{document}

\begin{xca}{6.4}
  Let $R$ be a ring.
  An element $x \in R$ is called \emph{nilpotent} if there exists a positive integer $n \in \mathbb{N}_{>0}$ such that $x^{n} = 0$.
  Show that the set of all nilpotent elements in $R$ is an ideal.

  You can try to \href{https://live.lean-lang.org/#codez=JYWwDg9gTgLgBAWQIYwBYBtgCMB0AlYAOwHMAVVAU2gE8cBJAEwqXRwCEkBnYAYwChQkWIhQZs%2BImUo0cAOWDpIMCoRjsuvPnwBuSKMCRZ0FOAAo8cAFxxS1MBQBUASjgBtAMIQQIAiTh4AXS1CJBAKTjAkHhNGZnQtJgAzOEIFKCQGXhYrOFjsi0sAXj4AbzgePX0KKCtCuBKADxyLAB9cznlFCGVVOAaAXz44OAAvaogAfTCQAHJauCxqIeG4AFpVuAAFKAgIZLQUOGB4HghVJCJOODQTMZ24CmMw1Rxl4c5oKCXhjIYpilm80WbzWG22u32qEOxyOVx46AgnAoDDgAFdCEwar9jsAzq8VnAPlAvstOCBUeh%2FoCigtvit1lsdntrlD4DDgHCEUiUejMYSKug9HByegYMAwJgKmK8SCiSTBnwVCi8uggA}{verify this exercise with Lean}.
\end{xca}

\subsection*{Hints for Lean}

\begin{leancode}
  import Mathlib.RingTheory.Ideal.Basic
  import Mathlib.RingTheory.Nilpotent.Basic

  variable (R : Type*) [CommRing R]

  namespace Ideal

  def nilradical : Ideal R :=
  { carrier := {x : R | IsNilpotent x}
    zero_mem' := by
      -- Proof that it contains the zero element.
      sorry
    add_mem' := by
      -- Proof that it is closed under addition.
      sorry
    smul_mem' := by
      -- Proof that it is closed under scalar multiplication.
      sorry
  }

  end Ideal
\end{leancode}

A \href{https://github.com/pedro-nunez/lean-algebra-2021/blob/master/nilradical/solution.lean}{solution} is available in the repository
\begin{center}
  \href{https://github.com/pedro-nunez/lean-algebra-2021}{https://github.com/pedro-nunez/lean-algebra-2021}.
\end{center}

{\renewcommand{\thefootnote}{\fnsymbol{footnote}}\footnotetext[2]{Exercises designed for the lecture \href{https://home.mathematik.uni-freiburg.de/arithgeom/lehre/ws21/azt}{\emph{Algebra and Number Theory}}, for which I was a teaching assistant.
This was an introductory lecture on abstract algebra, polynomial rings and Galois theory, taught by Annette Huber-Klawitter at the University of Freiburg during the Winter Semester 2021/2022.}}

\end{document}
